\chapter{Il gemello a campo intero del plug GENO}
\label{ch:genotwin}

Due programmi scritti in modo indipendente, in linguaggi diversi,
che si verificano a vicenda su un unico problema fisico sono un test
più forte di qualunque cosa ciascuno possa eseguire da solo: una
risposta condivisa non può venire da un errore condiviso. Questo
capitolo esegue quell'esperimento su un ugello a spina (aerospike)
--- uno la cui parete solida, la \emph{spina}, sta sotto lo scarico
mentre l'esterno del getto è delimitato da nient'altro che
l'atmosfera circostante. Il codice Fortran indipendente di progetto
GENO~\cite{geno} progetta una tale spina; il nostro solutore prende
quella forma come una comune parete assegnata, ricalcola l'intero
flusso attorno a essa, e i due campi di moto vengono confrontati
numero per numero. Il passo consegna la più forte verifica esterna
del programma --- e un risultato onesto.

\section{La domanda}

La sessione GENO a due vincoli (Capitolo~\ref{ch:status},
Sezione~\ref{sec:designinputs}) ha lasciato un residuo dichiarato:
il confronto a \emph{campo intero}. Non una manciata di grandezze
riassuntive, ma la nostra marcia del plug certificata eseguita sul
contorno di spina che il modulo RaoPlug di GENO ora produce in modo
pulito, confrontata campo contro campo --- pressione di parete
stazione per stazione, velocità nei punti interni di GENO, e massa e
spinta integrate. È la controparte per il plug della verifica fra
codici dell'ugello a campana del Capitolo~\ref{ch:referee}.

\section{Come funziona, per esteso}

\subsection{Che cosa occorre sapere prima}

Il solutore avanza un flusso supersonico lungo le due famiglie di
\emph{caratteristiche} --- le linee lungo cui una perturbazione può
viaggiare in flusso supersonico, inclinate dell'angolo di Mach da una
parte e dall'altra della direzione del flusso (Capitolo~\ref{ch:moc};
simboli nella Tabella~\ref{tab:notation}). Qui servono altre tre
parole.

Una \emph{linea iniziale} (o taglio di Cauchy, dal problema ai valori
iniziali che pone) è una linea tracciata attraverso il flusso sulla
quale lo stato completo --- posizione, velocità, pressione --- è
fornito come dato. Data una linea iniziale e le pareti, le equazioni
determinano tutto ciò che sta a valle di essa e nulla a monte. La
nostra è un taglio verticale attraverso il campo finito di GENO.

La \emph{linea dei valori iniziali} (IVL) è l'ultima caratteristica
della costruzione propria di GENO --- il bordo esterno della regione
che GENO calcola. Sotto di essa c'è il campo di progetto di GENO;
sopra, GENO non ha nulla da dire.

Il \emph{troncamento} compare qui in due sensi standard non
collegati, e il capitolo ha bisogno di entrambi: tagliare corta la
spina a una lunghezza finita (hardware), e l'errore che uno schema a
passo finito commette perché i suoi passi non sono infinitesimi
(numerica). Il risultato qui sotto riguarda il secondo.

\subsection{Due costruzioni indipendenti}

GENO costruisce la spina \emph{all'indietro}: ruota il ventaglio di
espansione al labbro del cowl a passi angolari fissi
($0.05^{\circ}$), integra una caratteristica per ogni membro del
ventaglio con le proprie celle, colloca ciascun punto di parete
richiedendo che la massa che scorre dentro quella caratteristica
eguagli la massa che la gola lascia passare (la terminazione di
massa, la N-65 di GENO) e taglia la spina nel punto $D$ in cui quel
bilancio è soddisfatto.

Il nostro motore non sa nulla di quella costruzione. È una marcia
assialsimmetrica generale --- celle interne, celle a parete assegnata
con la ricerca multi-colonna del piede, il bordo di getto libero,
ciascuna certificata nel Capitolo~\ref{ch:plug} --- e tratta il
contorno di GENO come una comune parete assegnata, marciando
\emph{in avanti} a partire da una linea iniziale. L'accordo è dunque
una prova sul problema continuo condiviso, non su codice condiviso.

\subsection{La posizione del problema: tutti i dati sono di GENO}

Non si assume nulla di analitico. Perfino il gas è identificato
\emph{dal campo di GENO}: il rapporto dei calori specifici
$\gamma = 5.34782609/4.34782609$ (la ricetta che il backend
termodinamico di GENO stesso implementa, un gas a calore specifico
costante), la costante del gas $R_g$ da $p/(\rho T)$ nei suoi punti
(dispersione $8.7\times10^{-9}$) e l'isentropica efficace del campo
stesso. Quest'ultima conta: il campo di GENO sta in realtà a una
pressione di camera $p_0 = 6.008172\times10^{6}$~Pa, \emph{non} al
nominale $6\times10^{6}$ --- uno scostamento uniforme della sua
costruzione, dispersione $1.1\times10^{-4}$ --- e porre il valore
nominale avrebbe fabbricato una differenza che non c'è.

La linea iniziale è un taglio verticale a $x_0 = 0.15$ attraverso le
caratteristiche di prima fase di GENO; ogni attraversamento è
interpolato sui punti propri di quella curva e verificato contro
un'interpolazione bidimensionale indipendente dello stesso campo, con
accordo a $10^{-5}$. La parete è il contorno di GENO, con pendenze
$\tan\theta_w$ misurate da esso (la sua tangenza è coerente a
$1.6\times10^{-5}$).

Sopra la linea dei valori iniziali giace una regione che GENO non
calcola, e la nostra marcia ha bisogno di qualcosa lì, quindi si pone
una striscia di getto libero, raccordata con continuità dallo stato
all'attraversamento della IVL allo stato sull'ultimo membro del
ventaglio. Non può contaminare il confronto --- un fatto sul flusso
supersonico, non una speranza, poiché ogni punto è influenzato solo
da ciò che sta a monte dentro il suo cono di dipendenza: ogni
caratteristica che arriva alla parete prima di $D$ lascia la linea
iniziale \emph{sotto} la IVL, e la IVL stessa atterra esattamente in
$D$, la fine della parete. È stato comunque verificato: sostituire
la striscia raccordata con una uniforme non cambia nulla sulla
parete.

\subsection{Soglie: dichiarate, con provenienza}

Ovunque altrove in questo programma la banda di superamento è
\emph{derivata}: lo stesso calcolo è eseguito a due risoluzioni, la
variazione misura l'errore proprio dello schema, e la banda è quella
variazione gonfiata da un fattore di sicurezza
(Capitolo~\ref{ch:method}). Qui i livelli sono invece
\emph{dichiarati} --- fissati in anticipo, con la loro provenienza
--- deliberatamente, per la ragione che è il risultato centrale di
questo passo qui sotto: la differenza misurata fra i due codici sta
\emph{sopra} ciò che il raffinamento di passo di ciascun codice
suggerisce, quindi una banda derivata dal raffinamento sarebbe il
metro sbagliato.

I livelli dichiarati vengono dall'accordo fra codici già stabilito
--- un integratore indipendente della curva di progetto, costruito
nella sessione a due vincoli, coincideva con GENO a tre o quattro
cifre significative --- rimisurato qui sul campo intero, con almeno
un fattore due di margine su ogni statistica: pressione di parete
$3\times10^{-3}$ (mediana misurata $1.9\times10^{-3}$), campo
assestato $q\ 10^{-3}$ / $\theta\ 1.5\times10^{-3}$ (95-esimo
percentile misurato $4.5\times10^{-4}$ / $5.3\times10^{-4}$),
integrali attraverso la frontiera condivisa $2.5\times10^{-3}$
(misurato $1.2\times10^{-3}$). Ciò che tiene onesti dei numeri
dichiarati sono i \emph{rigettatori} --- esecuzioni deliberatamente
corrotte che la suite deve far fallire (Capitolo~\ref{ch:method}):
dimostrano che lo strumento risolve difetti di un ordine di
grandezza superiori al divario che sta misurando. Una verifica
mantiene bande interamente derivate: l'identificazione del gas è
giudicata contro i pavimenti di autoconsistenza misurati di GENO
stesso (entalpia totale $2.0\times10^{-5}$, entropia
$4.7\times10^{-6}R$, pressione totale $1.1\times10^{-4}$).

\section{Che cosa è stato costruito}

\carrier{validation/a1\_geno\_plug\_twin.py}: l'estrattore di
riferimento (cartella di esecuzione GENO $\to$
\code{\_geno\_twin/ref.npz}: parete, linea dei valori iniziali, i
131{,}589 punti di campo di prima fase, identificazione del gas,
pavimenti), il costruttore del mondo della linea iniziale, la
quadratura di flusso calibrata --- lungo l'intera linea dei valori
iniziali di GENO riproduce la portata $\dot m$ stampata da GENO a
$2.4\times10^{-7}$ e la spinta $F$ a $2.3\times10^{-7}$, risolvendo
per strada una questione di definizione: la $F$ di GENO è
l'integrale di quantità di moto a pressione assoluta --- più le sei
verifiche di accordo G-1--G-6 e i due rigettatori R-1/R-2.

\section{Risultati di riferimento: 8/8 PASS}

G-1, il gas: le nostre tabelle riproducono pressione, temperatura e
densità di GENO nei suoi stessi punti a partire dalla sola velocità
locale, entro i pavimenti misurati di GENO (peggiore
$9.9\times10^{-5}$ in $p$). G-2: tutte le 25{,}034 celle certificate
da Newton a entrambe le risoluzioni. G-3, la parete: pressione entro
$3\times10^{-3}$ nel 100\% delle stazioni (mediana
$1.9\times10^{-3}$, peggiore $2.5\times10^{-3}$). G-4, l'interno:
superato il transitorio iniziale, il 97.7\% dei punti propri di GENO
concorda entro ($q\ 10^{-3}$, $\theta\ 1.5\times10^{-3}$); mediane
$3.8\times10^{-4}$ e $7.9\times10^{-5}$. G-5/G-6: massa e spinta
attraverso il tratto condiviso della linea dei valori iniziali
concordano a $1.2\times10^{-3}$ e $1.1\times10^{-3}$. Entrambi i
rigettatori scattano: un gas falsificato dell'$1\%$
($\gamma\times1.01$) fa crollare l'accordo di parete al 31\% delle
stazioni, un contorno scalato di $\times1.005$ lo fa crollare
allo 0\%.

\begin{figure}[htbp]
\centering
\includegraphics[width=\textwidth]{fig_geno_twin.pdf}
\caption{Il gemello a campo intero del plug GENO. (a)~Pressione di
parete lungo la spina: coordinata assiale $x$ [m] in orizzontale,
pressione sulla superficie della spina [bar] in verticale. La parete
propria di GENO (continua) e la nostra marcia in avanti sullo stesso
contorno (tratteggiata) si sovrappongono; la deviazione relativa
(punteggiata, asse destro, in unità di $10^{-3}$) resta sotto il
$3\times10^{-3}$ dichiarato in ogni stazione --- due codici, un
contorno, tre o quattro cifre. (b)~Lo stesso confronto dentro il
flusso: nei punti di campo propri di GENO, la differenza di modulo
della velocità fra i due codici, $|\Delta q|$ [m/s], contro $x$ [m]
(mediana e 95-esimo percentile per banda). Si guardi l'estremo
sinistro: il transitorio iniziale --- la nostra marcia che ingerisce
una linea iniziale tracciata attraverso la forte curvatura del flusso
vicino al labbro --- si esaurisce entro $x\approx0.35$~m, lasciando
un divario assestato di circa $4\times10^{-4}$ della velocità locale.
(c)~Che cosa sia quel divario assestato: la lunghezza effettivamente
marciata [m], variata spostando la linea iniziale, contro la
differenza di velocità che ne risulta alla parete $\Delta q$ [m/s]. I
punti cadono su una retta per l'origine ($\approx4\times10^{-4}$ di
$q$ per unità di lunghezza), mentre dimezzare i passi propri di GENO
sposta la parete di GENO di soli $3.5\times10^{-6}$ --- quindi non è
un errore locale in nessuno dei due codici ma una lenta deriva che si
accumula con la distanza marciata fra due solutori individualmente
convergenti, pari allo $0.5\%$ del contributo proprio del termine
sorgente assialsimmetrico.}
\label{fig:genotwin}
\end{figure}

\section{Il risultato: un divario fra codici misurato e con
attribuzione aperta}
\label{sec:twingap}

Il gemello certifica un accordo di tre-quattro cifre --- e la
differenza residua è essa stessa un risultato caratterizzato, non un
avanzo. Ha due parti. Vicino alla linea iniziale, un transitorio: a
quote intermedie l'angolo del flusso differisce fino a
$1.1\times10^{-2}$~rad, decadendo entro $x \approx 0.35$, invariato
al cambiare della risoluzione e mobile allo spostarsi della linea
iniziale --- la firma dell'ingestione di una linea iniziale
tracciata attraverso la forte curvatura del flusso vicino al labbro.
Si esaurisce. Ciò che resta è una \emph{deriva uniforme} assestata:
la nostra velocità corre $+4\times10^{-4}$ relativi più alta
sull'intera colonna, e quell'eccesso cresce linearmente con la
distanza marciata (Figura~\ref{fig:genotwin}c).

L'indagine ha escluso, per misura, ogni causa accessibile.
\emph{Il gas}: i due lati sono analiticamente equivalenti a circa
$10^{-8}$. \emph{La geometria della parete e la sua tangenza}:
$1.6\times10^{-5}$. \emph{I dati della linea iniziale}: concordi a
$10^{-5}$ per due vie di interpolazione indipendenti; raffinare la
risoluzione della poligonale da $10^{-4}$ a $5\times10^{-6}$ lascia
la deriva invariata. \emph{La striscia sopra la linea dei valori
iniziali}: sostituire il raccordo con uno stato uniforme non cambia
nulla, e il cono di dipendenza le vieta comunque di raggiungere la
parete. \emph{La discretizzazione di GENO}: una compilazione
diagnostica di GENO con il passo del ventaglio \emph{e} il passo
della curva entrambi dimezzati ha spostato la parete di GENO di
$3.5\times10^{-6}$, quindi GENO è convergente in passo (i sorgenti
sono poi stati ripristinati e il binario restaurato, con verifica
md5). \emph{La nostra risoluzione}: dimezzare ciascuno dei nostri due
conteggi di passo sposta la deviazione di meno del 10\% di sé
stessa.

Entrambi i campi soddisfano le stesse relazioni trapezoidali di
marcia a $10^{-8}$--$10^{-9}$ (un audit dei residui eseguito su
\emph{entrambe} le maglie). Ciò dimostra che i due sistemi discreti
sono algebricamente equivalenti --- e dimostra anche che un tale
audit è \emph{tautologico} come verifica di accuratezza: ogni codice
risolve esattamente le proprie relazioni, quindi un audit dei residui
non può mai vedere una distorsione che si accumula lentamente. Vale
la pena dirlo chiaramente, perché sembra una verifica forte e non lo
è. Infine, una marcia di controllo con il termine sorgente
assialsimmetrico spento dimensiona la deriva: è lo $0.5\%$ di ciò che
quel termine contribuisce alla velocità di parete (che è a sua volta
il $-5.4\%$ di $q$).

L'attribuzione resta onestamente \emph{aperta}. I meccanismi
candidati vivono precisamente dove il raffinamento non varia nulla
--- dal lato di GENO la contabilità della re-interpolazione durante
la sua marcia all'indietro, dal nostro l'interpolazione lungo la
corda trovata dalla ricerca del piede --- e dall'esterno, a questo
livello, non sono distinguibili. La domanda di competenza del
proprietario è depositata presso GENO. Per questo programma la
conseguenza operativa è quella già adottata: i confronti fra codici
con GENO portano il livello di accordo misurato di tre-quattro cifre
--- molto sopra l'errore di troncamento di entrambi i codici, molto
sotto qualunque cosa cui i gradienti dell'ottimizzatore siano
sensibili --- e i rigettatori tengono onesto lo strumento a quel
livello.
